\documentclass{article}
\usepackage{amsmath,amssymb,hyperref}
\title{Minimum notice for a handoff requiring two live permissions}
\date{26 November 2026}
\begin{document}
\maketitle
\paragraph{Finding.}
Two truthful reports with identical payload and arrival time can differ
in whether a local crew can finish a lawful handoff. For one recipient
view the exact required alphabet is the interval stabbing number.
For multiple private recipient views, with one common encoder that
does not know which view is present, the exact answer can already
be general graph coloring. Every view of the example individually
needs at most two symbols, but the common encoder needs three.
This is a fixed one-decision R5 model, not a field experiment or an
operational claim against an equally informed conditional actuator.

\paragraph{Timed authority and action.}
In a civil inspection task, management of an old load's support
moves from an incumbent organization to a receiving one. Each has a
separately narrowed permission that must be recognized by the actual
gate. Before physical clearance or after expiry, an available message
does not make the act lawful. Model old recognized permission on
integer ticks \([a,e)\), new on \([b,f)\), where these are the validity
and recognition actually enjoyed by the scoped actuator, not issuance
times printed on a document. New and old permissions are needed
throughout a \(h>0\)-tick stationary handshake; only the new one
is needed for the remaining \(r>0\) ticks of delivery.
Ready physical clearance from tick \(c\) stays true for the
tested motion, and no other rival job consumes the same bay or
support during it. A running load on a disjoint path continues
independently. A single irreversible start must be chosen by
decision cut \(d\), with start not earlier than \(d\), and delivery
due by end boundary \(D\). Finish on the last permitted tick
is processed before expiry at its next boundary. If an application
orders the expiry event before the simultaneous effect, this kernel
must change. Authenticated, unchanging names and no concurrent role
amendment within each hidden history are premises. Missing or late
recognition can make that history infeasible.

For a complete history the exact set of starts is the integer
interval
\[
 \big[\max(d,c,a,b),\ \min(e-h,\ f-h-r,\ D-h-r)\big].
\]
A start at the upper endpoint is allowed at the pinned boundary
ordering. The replay checks actual state changes in an authored
discrete plant, with the initially held load moving at handshake
completion and subsequently arriving at destination. A failure
late in the latter phase can leave the load already transferred;
it is not equivalent to refusal before starting. The independent
load follows its own disjoint cells and actually moves at ticks lying
inside each selected handoff, with a separate arrival at tick 10. No emergency rollback
or fresh reserve is presumed. The proof rule requires service
by the deadline, not just a successful handshake.
The intersection of scoped validity is an established idea
(e.g. experimental \href{https://datatracker.ietf.org/doc/html/rfc2693\#section-6.3}{SPKI
Certificate Theory, Section 6.3}); that reference does not supply
this physical timing rule or qualify sensor/authority data.

\paragraph{Two local views and the minimum notice.}
A sender observes key \(z\) and the recipient observes \(v\) by the
decision cut. Only the observable key may select the symbol, not the
complete history. The source law supplies their jointly possible
histories and the resulting actual feasible intervals, independently
of the sender encoding. The exact reply is not permitted to carry
authority or physical details outside a predeclared set of
timing notices; all symbol choices have equal slot/delivery budget.
No reversible probe, alternate route or late rechoice is in the
admitted action set. For each possible pair \((z,v)\), first
intersect the intervals of all histories with that pair to get \(J_{zv}\).
If this is empty, full disclosure of \(z\) is already insufficient
for \(v\), so no coding can repair the two observation interfaces.
An empty primitive interval separately rules out any start even
with the true state. If all intersections are nonempty, build graph
vertices from sender keys. Link \(z,z'\) whenever a common possible
view \(v\) has disjoint \(J_{zv},J_{z'v}\).

A correct code must give adjacent keys different symbols. For
sufficiency, color the graph. At every actually possible pair
of a color and a recipient view \(v\), all candidate intervals
for that view intersect pairwise; in one dimension their largest
left endpoint is no larger than their smallest right endpoint.
Choose that left endpoint. The resulting lookup function only
consumes \((v,\text{color})\) and produces a lawful start in all
compatible histories. The exact minimum symbol count is
\(\chi(G)\), or \(\lceil\log_2\chi(G)\rceil\) bits for a fixed-length
representation, excluding header and time acquisition costs.
This is a timed allowed-output specialization of the role of
confusability graphs in
\href{https://ieeexplore.ieee.org/document/1055607/}{Witsenhausen's
zero-error side information problem}, not a new general coding theorem.
If the feasible sets were not intervals, pairwise intersections
alone could fail; three allowed sets \(\{0,1\},\{1,2\},\{0,2\}\)
are a control. A varying clock discipline can itself invalidate the
interval decoder's model.

For one fixed \(v\), each symbol stands for one possible start,
so the answer equals the least number of times stabbing the
nonempty intervals. Take the earliest unhit right endpoint and repeat;
this hits every interval, while the witnesses for successive picks
are disjoint, giving the matching lower bound. This is the familiar
interval transversal argument (see
\href{https://www.renyi.hu/~tardos/2int.pdf}{Tardos}).
With windows \([2,4],[4,6],[7,7]\), a sender distinguishing the
third history can use one bit for starts 4 and 7. A sender
distinguishing the middle one instead cannot help: its other
key combines the first and third intervals. The same bit
received after the irrevocable choice also cannot help on this view,
which provides no common start by itself.

\paragraph{A strict penalty for private side information.}
Every simple graph can occur in this authored timing family.
For each edge give the recipient a distinct view possible with
just those two sender keys; use singleton starts 2 and 4 for
the two histories. These are actual instances of the timing
rule with a common \(d,h,r\), delayed clearance, and expiries
or deadlines yielding the chosen singleton. Include a private
view for each vertex, to allow even isolated keys.
There are no incompatibilities other than the designated edges.
In a pentagon, every individual view with an edge can use
two symbols; a common encoding requires three because the
graph is not bipartite, and a three-coloring exists.
Consequently there is no general greedy interval rule for the
common encoder. This result is about the permitted source relation
between two local views, not a claim that all these relations
hold in a particular multi-owner plant.

\paragraph{Comparison and audit.}
Explicit matching of every sender key grouping against all real
ticked starts and views is the direct reference; it agrees with
the exact construction. A direct qualified timing gate with the
whole necessary local observation before commitment can avoid
the remote notice. Such information or control cannot be assumed
when interpreting the one-notice bound. An ordinary conditional
authorization implementation with identical installed powers,
source clocks and gates has equivalent feasibility. No gain in
trust or actuating capability is claimed. The rule does not certify
that the gate sees actual clearance or correctly recognizes two
independent authorities in a deployed installation.
Scope is separate from the earlier
\href{https://github.com/eternal-november/eos-research-programmes/pull/621}{bay
consent width} calculation, which allowed one coarsely scoped query
and a response about a private assent vector, not this common
encoding of permissible times with recipient side information.

Replayed 24,576 discrete timed checks, all 32,767 nonempty interval
families on five ticks, and all 1,099 labeled simple graphs with
one through five vertices, 1,200 overlapping history tables, plus prepared controls for inadequate,
late and correctly keyed notices and a retained partial effect.
Three Lean primitives check exact half-open timing and the common
interval lemma. Compiled cases, roles and interval evidence are
authored and not independently selected. Broader R5, interface
validation, architectural advantage and acceptance remain open.
\end{document}
